\section{Resumen ejecutivo}

Este proyecto estudia qué características musicales y contextuales están asociadas con la
popularidad observada de canciones en Spotify. Su propósito es aportar evidencia para artistas,
sellos, marketing, curadores y gestión de catálogo, sin presentar las asociaciones como causas ni
como una receta de éxito. El trabajo sigue CRISP-DM y conserva trazabilidad mediante Git, notebooks
ejecutables, funciones en Python, pruebas automatizadas y documentación de progreso.

La fuente entregada por la asignatura contiene 114.000 filas y 21 columnas. No representa 114.000
canciones diferentes: solo existen 89.741 valores únicos de \code{track_id}; 16.641 IDs aparecen
más de una vez y generan 24.259 apariciones adicionales. Se detectaron tres celdas nulas, todas en
una fila que además posee duración cero, y una columna técnica \code{Unnamed: 0} que reproduce el
índice exportado. La distribución por género original contiene exactamente 1.000 filas para cada
una de 114 etiquetas, indicio de un diseño de muestra balanceado y no de la prevalencia real del
catálogo.

La preparación retiró la columna técnica y la única fila inválida. Después consolidó cada
\code{track_id}, verificando que sus atributos musicales y textuales fueran invariantes. Las
etiquetas de género se conservaron como una lista ordenada separada por barras verticales en
\code{track_genres}; cuando un ID presentaba diferentes valores históricos de popularidad, se usó
la mediana. El resultado contiene 89.740 canciones únicas, 20 columnas, cero nulos, cero IDs
duplicados y duración estrictamente positiva.

\begin{table}[H]
\centering
\caption{Comparación reproducible del dataset antes y después de la preparación.}
\label{tab:comparacion-resumen}
\resizebox{\textwidth}{!}{\begin{tabular}{lrrrrrr}
\toprule
 & Filas & Canciones unicas & Media & Mediana & Porcentaje igual a 0 & Porcentaje mayor o igual a 70 \\
dataset &  &  &  &  &  &  \\
\midrule
Original & 114000 & 89741 & 33.239 & 35.000 & 14.053 & 4.800 \\
Procesado & 89740 & 89740 & 33.198 & 33.000 & 10.468 & 3.480 \\
\bottomrule
\end{tabular}
}
\end{table}

El EDA procesado obtiene una popularidad media de 33,20, mediana de 33 y desviación estándar de
20,58. El 10,47\% de las canciones tiene popularidad cero y el 3,48\% alcanza 70 o más. Ningún
atributo acústico individual muestra una asociación fuerte con el objetivo. La mayor magnitud es
\code{instrumentalness} (Pearson $-0{,}128$; Spearman $-0{,}124$); \code{loudness} y
\code{danceability} presentan asociaciones positivas pequeñas. Entre predictores sí aparecen
redundancias relevantes: \code{energy} con \code{loudness} ($0{,}759$) y con
\code{acousticness} ($-0{,}733$).

Las canciones marcadas como explícitas tienen una media 4,03 puntos superior y una mediana cuatro
puntos superior, pero el tamaño de efecto estandarizado es pequeño ($d=0{,}196$). Esto distingue
una diferencia visible en una muestra grande de una diferencia práctica considerable. El ranking
por género es estable al exigir soportes mínimos de 100, 200 o 500 canciones, aunque esa estabilidad
no corrige el muestreo artificialmente balanceado ni convierte las diferencias en efectos causales.

Los riesgos principales son la falta de temporalidad, variables contextuales omitidas, posible
memorización de artistas, leakage mediante identificadores, alta cardinalidad y deriva del concepto
de popularidad. Éticamente, un modelo podría reforzar ventajas existentes, reducir diversidad o
automatizar promoción de forma injustificada. El dataset no contiene datos directos de oyentes, por
lo que el riesgo de privacidad individual es relativamente bajo; aumentaría si se integraran
perfiles, ubicaciones o historiales de escucha.

El proyecto está preparado para una etapa inicial de modelamiento de regresión. Se recomienda
construir una línea base, excluir \code{track_id}, codificar géneros como multi-hot, evaluar con
MAE, RMSE y $R^2$, y documentar desempeño por segmentos. El resultado debe utilizarse como apoyo a
decisión humana, no como criterio único de mérito artístico, inversión o promoción.
